\scanpage{7}
\mspage{6}\hfill\mspage{3}
pour tout diagramme
\diagram{\[
\left\{\begin{gathered}\begin{tikzcd}[row sep=large]
\deleted{A}\inserted{X}&\\
{|X|}\arrow[r,"{\substack{\alpha\\\sim}}"]&E
\end{tikzcd}\end{gathered}\right.
\]}
avec $X\in\mathrm{Ob}\,\mathrm{Esp}_{\mathcal M_*}$, $\exists$ une unique \deleted{$Y$ et} flèche $f:X\mathbin{-}Y$ de $\mathrm{Esp}_{\mathcal M_*}$, qui relève $\alpha$ (``Transport de structure'').

\uline{Proposition} Dans la catégorie $\mathrm{Esp}_{\mathcal M_*}$ des espaces $\mathcal M_*$-modérés, il y a des
\deleted{produits fibrés \inserted{\deleted{(de deux objets} sur un troisième)}, et des produits (de deux objets), et le foncteur $X\mathrel{\mapstochar\relbar}|X|$ y commute. Si on suppose
\[
(M5)\qquad R_0\ne\varnothing
\]
alors il \deleted{foncteur} y a aussi un objet final (qui est uniponctuel) et le foncteur}
\inserted{des limites finies}
et le foncteur $X\mathrel{\mapstochar\relbar}|X|$ y commute.
\margin{\unsure{trivialité de} \unsure{suivi} d'un \illegible{1} \unsure{et des} \illegible{2} des \unsure{foncteurs}}

\uline{Proposition} Si pour \uline{un} $n\in\mathbb N^*$ on a \deleted{$\unsure{R}$} $R^n\in\mathcal M_n$, alors c'est vrai pour $\forall n$, et même pour $\forall I$ fini \inserted{dans}
\[
R^I\in\mathcal M_I.
\]
Considérons l'axiome (\uline{facultatif})
\[
\deleted{\unsure{M5}}(M0)\qquad R\in\mathcal M_1\qquad(\Longrightarrow\forall I\in(\mathrm{Ens\ fini}),\ R^I\in\mathcal M_I)
\]
\margin{Quand on n'a pas $(M0)$, mais que l'on suppose $(M5)$, alors la réunion des $A\in\mathcal M_1$, fait que $\unsure{R'}$ est \unsure{obtenu} \unsure{comme} \illegible{2} \unsure{de} $\unsure{R}$ \unsure{réunion} \unsure{filtrante} \illegible{2} $\unsure{\Sigma}$.\par
\unsure{On dispose} de \unsure{ces} \unsure{ens.} \unsure{stables} \illegible{2} l'axiome $(M1)$ \unsure{s'exprime} \illegible{1} \unsure{l'influence} \unsure{de} \unsure{foncteurs} \illegible{2}.\par
\unsure{comme} $(M5)$ l'expression \unsure{espace} \unsure{muni de} \unsure{l'} $\unsure{R}$ \unsure{aussi} \unsure{vide} \unsure{annulé}.}
\margin{\unsure{objet} \unsure{abs.} \unsure{Cat}\uline{\unsure{f}}}
